\documentclass[a4paper]{article}

\usepackage{amssymb} 
\usepackage[a4paper]{geometry}
\usepackage[T1]{fontenc}
\usepackage[utf8]{inputenc}
\usepackage{graphicx}
\graphicspath{ {images/} }
\usepackage{todonotes}
\usepackage{hyperref}
\hypersetup{
    colorlinks=true,
    linkcolor=black,
    filecolor=magenta,
    urlcolor=blue
}
% changes font
\usepackage[style=verbose-ibid,backend=bibtex]{biblatex}
\bibliography{wipro-main-doc}

\usepackage{times}
% Set name of image label
\renewcommand{\figurename}{Abbildung}

\title{
    {Wirtschaftsprojekt} \\
    \vspace{10mm}
    { Herbstsemester 2022 } \\
    \vspace{10mm}
    % https://commons.wikimedia.org/wiki/File:HSLU_2022_logo.svg
    {\includegraphics[width=75mm]{img/hsluLogo2022.png}}
}

\author{Yannis Krämer und Nicolas Wiedmer}
% TODO: Update date
\date{22.09.2022}

\begin{document}

\maketitle

\newpage

\noindent
\fontsize{12}{14}
\textbf{Wirtschaftsprojekt an der Hochschule Luzern -- Informatik} \\ \vspace*{0.6cm}

\fontsize{10.95}{12}
\noindent
\textbf{Titel:} STAIR Discord Bot \\ \vspace*{0.2cm}

\noindent
\textbf{Studentin/Student:} Yannis Krämer \newline \newline
\textbf{Studentin/Student:} Nicolas Wiedmer \newline \newline
\textbf{Studiengang:} BSc Informatik oder Wirtschaftsinformatik  \newline \newline
\textbf{Jahr:} 2022 \newline \newline
\textbf{Betreuungsperson:} Markus Waldmann \newline \newline
\textbf{Expertin/Experte:} \newline \newline
\textbf{Auftraggeberin/Auftraggeber:} STAIR (Martin Steiger \& Estefania Otero)\newline \newline \newline
\textbf{Codierung / Klassifizierung der Arbeit:}\\
$\boxtimes$ \"Offentlich 
$\square$ Vertraulich


%%% you can use \boxtimes for filling a cross inside the square
%%% e.g., $\boxtimes$ A: Einsicht 	(Normalfall) 


\paragraph{\textbf{Eidesstattliche Erkl\"arung}}
Ich erkl\"are hiermit, dass ich/wir die vorliegende Arbeit selbst\"andig und ohne unerlaubte fremde Hilfe angefertigt haben, alle verwendeten Quellen, Literatur und andere Hilfsmittel angegeben haben, w\"ortlich oder inhaltlich entnommene Stellen als solche kenntlich gemacht haben, das Vertraulichkeitsinteresse des Auftraggebers wahren und die Urheberrechtsbestimmungen der Hochschule Luzern respektieren werden. \newline \newline
Ort / Datum, Unterschrift	\underline{\hspace*{8cm}} \newline \newline
Ort / Datum, Unterschrift	\underline{\hspace*{8cm}} \newline \newline \newline
\textbf{Ausschliesslich bei Abgabe in gedruckter Form: \\
Eingangsvisum durch das Sekretariat auszuf\"ullen}\newline \newline
Rotkreuz, den	\underline{\hspace*{4cm}}	\hspace*{1cm} Visum:	\underline{\hspace*{4cm}}

\newpage
\section*{I{\hspace*{1cm}}Abstract}
\todo{schreiben}

\newpage

\tableofcontents
\newpage
\section{Problem, Fragestellung, Vision}

\todo{schreiben}

\section{Stand der Technik}

\subsection{Sprach- und Text-Messenger}

\subsection{Discord}
Der Onlinedienst Discord ist ein Instant-Messaging und Chat Tool mit dem auch Sprach- und Videokonferenzen 
abgehalten werden können. Er kann online auf einer Webseite aufgerufen werden, oder mithilfe eines Clients lokal. 
Discord unterstützt alle gängigen Betriebssysteme und kann auch auf mobilen Endgeräten verwendet werden.

\subsubsection{Geschichte}
Ursprünglich wurde Discord für Computerspiele geschaffen, um nebenher mit seinen Spielkameraden zu texten oder
miteinander zu reden. Das Ziel war es, neben der komplexen Spielmechanik einen Chat-Messanger zu bauen, der 
benutzerfreundlich und effizient ist. 2012 wurde das Unternehmen Discord Inc (damals noch Hammer \& Chisel)
als Startup gegründet.\autocite{} 2014 erhielt Hammer \& Chisel für die Weiterentwicklung von Discord zusätzliche
Finanzierungsmittel von verschiedenen Unternehmen.

2015 wurde Discord unter der Domain discordapp.com veröffentlicht. Der Text und Sprach Messanger wurde immer beliebter
und erhielt immer mehr Finanzierungskosten, wie 2016 20 Millionen von WarnerMedia \autocite{} .
2018 kündete auch Microsoft Discord Unterstützung für ihre XBox Live Nutzer an und unterstützte Discord mit einer Finanzierung
von 150 Millionen US-Dollar. Bewertet wurde Discord Inc nun auf etwa 2 Milliarden US-Dollar.

Aufgrund der Covid-19 Pandemie und den steigenden Benutzerzahlen stellte sich Discord auch für andere Zielgruppen ein. 
So wurde der Fokus von Videospielen weggelenkt auf einen universellen Kommunikations-und Chat Client, um es für Branchen 
wie Schule oder Arbeit ansprechender zu machen. Damit änderte das Unternehmen ihr Motto von \textit{Chat for Gamers} zu 
\textit{Chat for Communities and Friends} und führte Servervorlagen ein.

Heute hat Discord mehr als 140 Millionen monatlich aktive User, verwaltet etwa 13.5 Millionen aktive Server und
wird auf 15 Milliarden US-Dollar gewertet. (Alles Stand 2021) \autocite{}

\autocite{discord_wiki}


\subsection{Discord Einschränkungen}

% https://netcord.site/new-discord-channel-creation-ratelimit/
% https://support.discord.com/hc/en-us/community/posts/360056762431-Increase-channel-limit
% https://support.discord.com/hc/en-us/community/posts/360032363631-Increase-the-Max-Number-of-Roles-per-Server

\section{Ideen und Konzepte}

\todo{einfügen}

\section{Methoden}

\todo{einfügen}

\subsection{Technologien}

\todo{einfügen}
% https://www.nuget.org/packages/linq2db.MySql/
% http://www.primaryobjects.com/2009/01/24/using-mysql-and-linq-to-sql-in-c-asp-net/

\section{Realisierung}

\todo{schreiben}

\section{Evaluation und Validation}

\todo{schreiben}

\section{Ausblick}

\todo{schreiben}

\newpage

\section{Anhänge}

\todo{schreiben}

\section{Glossar}

% https://blog.aristolo.com/de/definition-begriffe-bachelorarbeit-masterarbeit-dissertation/
\todo{schreiben}

\section{Abbildungsverzeichnis}
\listoffigures

\section{Tabellenverzeichnis}
\listoftables

\section{Literaturverzeichnis}
\printbibliography

\end{document}
